\subsection{淡收敛的一个性质}

引理 5. --- 设 X 是一个局部紧空间，\(\mu\) 是 X 上的一个有界正测度，F 是一个 Banach 空间，f 是 X 上的一个取值于 F 的有界函数。下列条件等价：

\begin{enumerate}
\def\labelenumi{(\roman{enumi})}
\item
  \(\mathbf{f}\) 的不连续点所成之集是 \(\mu\)-可忽略的。
\item
  对每个 \(\varepsilon > 0\)，存在元素 \(\mathbf{a}_1, \ldots, \mathbf{a}_n\)（\(F\) 的）、函数 \(g_1, \ldots, g_n\)（属于 \(\mathcal{K}(X)\)），以及一个有界连续函数 \(h \geqslant 0\)（\(X\) 上的），使得 \(|\mathbf{f} - g_1 \mathbf{a}_1 - \cdots - g_n \mathbf{a}_n| \leqslant h \leqslant 2 \sup |\mathbf{f}|\) 在 \(X\) 上成立，且 \(\int h d\mu \leqslant \varepsilon\) 。用 \(N\) 表示 \(\mathbf{f}\) 的不连续点所成之集，并令 \(M = \sup |\mathbf{f}|\) 。
\setcounter{enumi}{0}
\item
  \(\Rightarrow\) (ii)。设条件 (i) 满足。设 \(\varepsilon > 0\) 。函数 f 是 \(\mu\)-可积的（No.~2, 命题 5 的推论 4 与 No.~6, 定理 5），因此存在 \(a_{1}, \ldots, a_{n}\)（F 中的）与 \(g_{1}, \ldots, g_{n}\)（\(\mathcal{K}(X)\) 中的），使得在令 \(k = |f - g_{1}a_{1} - \cdots - g_{n}a_{n}|\) 时，有 \(\int k d\mu \leqslant \varepsilon/2\)（§3, No.~5,
命题 10）。将 \(g_1, \ldots, g_n\) 乘以 \(\mathcal{K}(X)\) 中一个适当的同一元素，我们还可以假设
\end{enumerate}

\[
\left| g _ {1} \mathbf {a} _ {1} + \dots + g _ {n} \mathbf {a} _ {n} \right| \leqslant \mathrm{M} = \sup | \mathbf {f} |
\]

在 X 上成立，由此 \(k \leqslant 2M\) 。集 \(N'\)——\(k\) 的不连续点所成——含于 \(N\)，故可忽略。对每个 \(x \in X\)，令 \(l(x) = \limsup_{y \to x} k(y)\) 。

则在 X 上 \(2\mathrm{M} \geqslant l \geqslant k\)，且 \(l = k\) 在 \(\mathrm{X} - \mathrm{N}'\) 上成立（即关于 \(\mu\) 几乎处处），因此 \(\int l d\mu \leqslant \varepsilon / 2\) 。另一方面，\(l\) 有界且上半连续，故是 \(\geqslant l\) 的有界连续函数所成之集的下包络。因此存在 X 上的有界连续函数 \(h \geqslant l\)，使 \(h \leqslant 2\mathrm{M}\) 且 \(\int h d\mu \leqslant \int l d\mu + \varepsilon / 2\)（§4, No.~4, 命题 5 的推论 2）。于是 \(\int h d\mu \leqslant \varepsilon\)，且 \(|\mathbf{f} - g_1\mathbf{a}_1 - \dots - g_n\mathbf{a}_n| \leqslant h\) 。

\begin{enumerate}
\def\labelenumi{(\roman{enumi})}
\setcounter{enumi}{1}
\tightlist
\item
  \(\Rightarrow\) (i)。设条件 (ii) 满足。对每个 \(x \in X\)，用 \(\omega(x)\) 表示 f 在 x 处的振幅（GT, IX, §2, No.~3）。设 \(\varepsilon > 0\) 。存在具有 (ii) 中性质的 \(a_{1}, \ldots, a_{n}, g_{1}, \ldots, g_{n}, h\)。对每个 \(x \in X\)，\(\omega(x)\) 是 \(f - g_{1}a_{1} - \cdots - g_{n}a_{n}\) 在 x 处的振幅，因此 \(\omega(x) \leqslant 2h(x)\) 。于是 \(\int \omega d\mu \leqslant 2\varepsilon\) 。从而，集 \(A_{\varepsilon}\)——由 \(x \in X\)（满足 \(\omega(x) \geqslant \sqrt{\varepsilon}\) 者）组成——适合 \(\mu(A_{\varepsilon}) \leqslant 2\sqrt{\varepsilon}\) 。这就证明了 \(\mu(N) \leqslant 2\sqrt{\varepsilon}\)，故 \(\mu(N) = 0\) 。
\end{enumerate}

命题 22. --- 设 F 是一个 Banach 空间，X 是一个局部紧空间，E 是 X 上有界正测度所成之集，μ 是 E 的一个元素，而 B 是 E 上的一个滤子基。假设 B 淡收敛于 μ，且 \textbar\textbar ν\textbar\textbar{} 关于 B 趋于 \textbar\textbar μ\textbar\textbar{}。设 f 是 X 到 F 中的一个映射，满足下列条件：

\begin{enumerate}
\def\labelenumi{(\roman{enumi})}
\item
  \(\mathbf{f}\) 有界，且关于 \(\mu\) 以及关于属于 \(\mathfrak{B}\) 的某个元素的每个测度都可积；
\item
  \(\mathbf{f}\) 的不连续点所成之集是 \(\mu\)-可忽略的。
\end{enumerate}

则 \(\int \mathbf{f}d\nu\) 趋于 \(\int \mathbf{f}d\mu\)（关于 \(\mathfrak{B}\)）。

设 \(\varepsilon > 0\) 。存在元素 \(\mathbf{a}_1, \ldots, \mathbf{a}_n\)（\(F\) 中的）、函数 \(g_1, \ldots, g_n\)（\(\mathcal{K}(X)\) 中的），以及一个有界连续函数 \(h \geqslant 0\)（\(X\) 上的），使得

\[
\left| \mathbf {f} - g _ {1} \mathbf {a} _ {1} - \dots - g _ {n} \mathbf {a} _ {n} \right| \leqslant h \leqslant 2 \sup | \mathbf {f} |
\]

在 X 上成立，且 \(\int h d\mu \leqslant \varepsilon\)（引理 5）。令 \(M = \sup |f|\) 。存在 X 的紧子集 K，使 \(\mu^{*}(X - K) \leqslant \varepsilon\)（§4, No.~7, 命题 12 与 No.~6, 定理 4）、K 在 X 中的一个紧邻域 \(K'\)，以及一个连续映射 \(h'\)（X 到 {[}0,2M{]} 中的），使 \(h' = h\) 在 K 上、\(h' = 2M\) 在 \(X - K'\) 上成立；把 \(h'\) 替换为 \(\sup(h, h')\)，我们还可以假设 \(h' \geqslant h\) 。则 \(\int(h' - h) d\mu \leqslant 2M\mu^{*}(X - K) \leqslant 2M\varepsilon\) 。另一方面，\(h' = h_1 + 2M\)，其中 \(h_1 \in \mathcal{K}(X)\) 。注意到 §4, No.~7 命题 12，数 \(\int h^{\prime}d\nu = \int h_{1}d\nu +2\mathrm{M}\| \nu \| \text{ 趋于 } \int h_{1}d\mu +2\mathrm{M}\| \mu \| = \int h^{\prime}d\mu \text{ 关于 } \mathfrak{B}.\) 于是存在一个 \(A\in \mathfrak{B}\)，使对每个 \(\nu \in A\)，

\[
\begin{array}{c} \Big | \int (g _ {1} \mathbf {a} _ {1} + \dots + g _ {n} \mathbf {a} _ {n}) d \nu - \int (g _ {1} \mathbf {a} _ {1} + \dots + g _ {n} \mathbf {a} _ {n}) d \mu \Big | \leqslant \varepsilon , \\ \int h d \nu \leqslant \int h ^ {\prime} d \nu \leqslant \int h ^ {\prime} d \mu + \varepsilon \leqslant \int h d \mu + 2 \mathrm{M} \varepsilon + \varepsilon \leqslant 2 (\mathrm{M} + 1) \varepsilon . \end{array}
\]

这些不等式蕴涵

\[
\begin{array}{c} \Big | \int \mathbf {f}   d \nu - \int \mathbf {f}   d \mu \Big | \leqslant \\ \int h   d \nu + \Big | \int (g _ {1} \mathbf {a} _ {1} + \dots + g _ {n} \mathbf {a} _ {n})   d \nu - \int (g _ {1} \mathbf {a} _ {1} + \dots + g _ {n} \mathbf {a} _ {n})   d \mu \Big | + \int h   d \mu \\ \leqslant 2 (\mathrm{M} + 2) \varepsilon , \end{array}
\]

这就证明了命题。

注. --- 命题 22 的条件 (i) 与 (ii) 当 f 连续且有界时满足。

例. --- 取 X 为模等于 1 的复数组成的紧空间 U。令 \(\mu(f) = \int_{0}^{1} f(e^{2i\pi t}) dt\)（对每个 \(f \in \mathcal{K}(U)\)），便定义了 U 上质量为 1 的一个正测度。另一方面，设 \(\theta\) 是一个实数；对每个整数 \(n \geqslant 0\)，设 \(\nu_{n}\) 是放置在 U 的点 \(e^{2i\pi n\theta}\) 处的单位质量，并令

\[
\mu_ {n} = \frac {1}{n + 1} (\nu_ {0} + \dots + \nu_ {n}),
\]

于是 \(\mu_n\) 是 \(\mathbf{U}\) 上质量为 1 的正测度。那么，若 \(\theta\) 是无理数，则 \(\mu_n\) 淡收敛于 \(\mu\) 。因为，由于函数 \(z \mapsto z^k\)（ \(k \in \mathbf{Z}\) ）的线性组合在 \(\mathcal{K}(\mathbf{U})\) 中稠密（GT, X, §4, No.~4, 命题 8），只需证明 \(\mu_n(z^k)\) 趋于 \(\mu(z^k)\)（对 \(k \in \mathbf{Z}\)） 。而对 \(k = 0\)，\(\mu_n(z^k) = \mu(z^k) = 1\)；对 \(k \neq 0\)，

\[
\mu_ {n} (z ^ {k}) = \frac {1}{n + 1} (1 + e ^ {2 i \pi k \theta} + e ^ {4 i \pi k \theta} + \dots + e ^ {2 i \pi k n \theta}).
\]

由于 \(e^{2i\pi k\theta} \neq 1\)（因为 \(\theta\) 是无理数），我们推出

\[
\left| \mu_ {n} (z ^ {k}) \right| = \left| \frac {1}{n + 1} \frac {e ^ {2 i \pi k (n + 1) \theta} - 1}{e ^ {2 i \pi k \theta} - 1} \right| \leqslant \frac {1}{n + 1} \frac {2}{\left| e ^ {2 i \pi k \theta} - 1 \right|},
\]

因此 \(\mu_n(z^k)\) 趋于 \(0 = \mu (z^k)\) 。在这些条件下，命题 22 可以应用，我们特别看出：若 A 是 U 的一个其边界关于 \(\mu\) 可忽略的子集，则 \(\mu_n(A)\) 趋于 \(\mu(A)\) 。换句话说，若用 \(p_n\) 表示整数 \(k \in [0, n]\)（满足 \(e^{2i\pi k\theta} \in A\) 者）的个数，则 \(n^{-1}p_n\) 趋于 \(\mu(A)\)（当 \(n\) 趋于 \(+\infty\) 时）。

